%----------------------------------------------------------------------------
\chapter{Glossary}\label{app:glossary}
%----------------------------------------------------------------------------
This chapter defines process terms (\cref{tab:glossary-process}) and technical terms (\cref{tab:glossary-technical}). It also lists English--Hungarian pairs for the Hungarian abstract and discussions of your work (\cref{tab:glossary-hu}).

\section{Process terms}
{\small
\begin{xltabular}{\linewidth}{>{\raggedright\arraybackslash}p{0.24\linewidth} X}
  \caption{Process terms.}\label{tab:glossary-process}\\
  \toprule \textbf{Term} & \textbf{Meaning} \\ \midrule \endfirsthead
  \toprule \textbf{Term} & \textbf{Meaning} \\ \midrule \endhead
  \bottomrule \endlastfoot
  Topic manifest & One page approved before the semester: the problem, why it matters, the expected contribution. \\
  Literature review & A search for sources, logged in one notes file with its search strings, databases and dates; there is no screening protocol. It gives the gap claim a documented set of sources. \\
  Gap analysis & Finding what existing approaches cannot do. Produces the positioning figure, gap matrix and a one-sentence gap. \\
  Gap matrix & Requirements compared against approaches in \tplref{ch:related}; see \cref{app:chapters:existing-approaches}. \\
  Positioning figure & Related work plotted on two axes with three levels each (\tplref{fig:positioning}); see \cref{app:chapters:existing-approaches}. \\
  Leading figure (vezérábra) & The one figure that explains the work; see \cref{app:writing:figure}. \\
  Spine & One page fixing the paper's claims, terminology and notation; see \cref{app:writing:spine}. \\
  Skeleton mind map & Each chapter planned as paragraph topic sentences with supporting evidence, before drafting prose; see \cref{app:writing:spine}. \\
  Dependency graph & The skeleton's paragraphs with an arrow to each paragraph that needs them, the contribution in green; see \cref{app:writing:depgraph}. \\
  Beats & What each planned paragraph establishes, its evidence, its use of the example and its link to the next paragraph. \\
  Essentials & The elements every paper of this kind contains, measured by the lint on every compile; see \cref{app:writing:essentials}. \\
  Premortem & Predicting why the paper might score poorly on each criterion, then addressing those objections before drafting. TDK papers receive placements rather than acceptance or rejection. \\
  Anchor & An artefact we approve before the next phase and use to check later work: the topic manifest and the spine, one page each, and the full draft, approved at a revision recorded in the consultation summary; see \cref{app:lifecycle}. \\
  Drift & Work that contradicts an anchor without re-opening it; see \cref{app:lifecycle}. \\
  Official step & A step of the faculty's TDK with a portal date: registration with the abstract, paper upload, the conference; see \cref{app:tdk:dates}. \\
  Department session & The writing session before the deadline and the presentation rehearsal before the conference, organised by the department for all its TDK students; see \cref{tab:lifecycle}. \\
  Context upload & Putting our papers, notes and earlier work on the topic into \texttt{notes/} of your notes repository. \\
  Hill climb & Repeatedly assessing work and fixing its weakest point, whether proofs, implementation or, in the review sheet hill climb, the paper against the review sheet; \cref{tab:lifecycle} says when each is done. \\
  Reviewer pass & An agent's reading as a simulated reviewer with a narrow brief, before real review; see \cref{app:lifecycle:reviewers}. \\
  Skim view & The reviewer PDF reduced to headings, labelled items, boxes, figures and captions, with paragraph lengths replacing prose; paper CI builds it on every push, see \cref{app:lifecycle:reviewers}. \\
  Requirements & Every requirement in the guide, with its source, its check and its due phase, including lint rules (\cref{app:code:paperci}); see \cref{tab:polish}. \\
  Burn-down & Counts of todos, beats, placeholder figures and TODO cells, which paper CI reports on every push and the lifecycle CI records (\cref{app:code:paperci}); \cref{tab:polish} says when they reach zero. \\
  Reviewer PDF & The paper built from \texttt{reviewer.tex}: no notes, no guide; the file to upload. \\
  Results tag & The Git tag and Zenodo version used for the paper's numbers; see \cref{app:zenodo:cite,app:zenodo:versions}. \\
  Version DOI, concept DOI & Identifiers for one archived release and for all versions together, respectively; \cref{app:zenodo:cite} says which one the paper cites. \\
\end{xltabular}
}

\section{Technical terms}
The common glossary, \texttt{include/glossary-common.tex}, declares the abbreviations, notation and terms our topics share (\cref{tab:glossary-common}). Use them with \texttt{\textbackslash gls\{key\}} without declaring them: only the ones your text uses appear in your lists. To change an entry for your paper, declare the same key in \texttt{content/glossary.tex}, and yours replaces it. The template's running example and the tools of our topics are not in the common glossary; \cref{tab:glossary-technical} explains them.

% The rows of tab:glossary-common come from the keys \commonentry listed (include/preamble.tex);
% \glsentry... commands print a field without marking the entry as used in the paper's lists.
\makeatletter
\newcommand*{\tdk@abbrrows}{}\newcommand*{\tdk@symrows}{}\newcommand*{\tdk@termrows}{}
\newcommand{\tdk@abbrrow}[1]{\Glsentrylong{#1} (\glsentryshort{#1}) & \texttt{#1} & \glsentryuseri{#1}. \\}
\newcommand{\tdk@symrow}[1]{\glsentryname{#1} & \texttt{#1} & \glsentrydesc{#1}. \\}
\newcommand{\tdk@termrow}[1]{\Glsentryname{#1} & \texttt{#1} & \glsentrydesc{#1}. \\}
\newcommand{\tdk@sortcommon}[1]{%
  \edef\tdk@type{\glsentrytype{#1}}%
  \ifdefstring{\tdk@type}{abbreviations}{\gappto\tdk@abbrrows{\tdk@abbrrow{#1}}}{%
    \ifdefstring{\tdk@type}{symbols}{\gappto\tdk@symrows{\tdk@symrow{#1}}}{\gappto\tdk@termrows{\tdk@termrow{#1}}}}}
\forlistloop{\tdk@sortcommon}{\tdkcommonkeys}
{\small
\begin{xltabular}{\linewidth}{>{\raggedright\arraybackslash}p{0.27\linewidth} l X}
  \caption{The common glossary, \texttt{include/glossary-common.tex}.}\label{tab:glossary-common}\\
  \toprule \textbf{Entry} & \textbf{Key} & \textbf{Meaning} \\ \midrule \endfirsthead
  \toprule \textbf{Entry} & \textbf{Key} & \textbf{Meaning} \\ \midrule \endhead
  \bottomrule \endlastfoot
  \multicolumn{3}{@{}l}{\textit{Abbreviations}} \\
  \tdk@abbrrows \midrule
  \multicolumn{3}{@{}l}{\textit{Notation}} \\
  \tdk@symrows \midrule
  \multicolumn{3}{@{}l}{\textit{Terms}} \\
  \tdk@termrows
\end{xltabular}
}
\makeatother

{\small
\begin{xltabular}{\linewidth}{>{\raggedright\arraybackslash}p{0.24\linewidth} X}
  \caption{Technical terms outside the common glossary.}\label{tab:glossary-technical}\\
  \toprule \textbf{Term} & \textbf{Meaning} \\ \midrule \endfirsthead
  \toprule \textbf{Term} & \textbf{Meaning} \\ \midrule \endhead
  \bottomrule \endlastfoot
  Verifiable credential (VC) & A digitally signed statement about a subject, issued by an issuer, held by a holder, and checked by a verifier (W3C standard). \\
  Zero-knowledge proof (ZKP) & A proof that a statement is true which reveals nothing beyond the statement itself. \\
  Hyperledger Fabric & A permissioned blockchain platform for consortia of known organisations. \\
  Tamarin & A prover for security protocols that searches for attacks under a symbolic adversary model. \\
  Lean & A proof assistant and programming language for machine-checked proofs. \\
  Refinery & A tool for partial modelling and graph generation over metamodels. \\
\end{xltabular}
}
VC and ZKP are the running example's abbreviations, declared in \texttt{content/glossary.tex}. Add the terms only your topic needs there too, and check their definitions against the topic's literature.

\section{English--Hungarian pairs}
Use the same Hungarian term for each concept throughout the abstract and presentation. These pairs come from earlier TDK abstracts, including variants where they were used.
{\small
\begin{xltabular}{\linewidth}{>{\raggedright\arraybackslash}p{0.36\linewidth} >{\raggedright\arraybackslash}X}
  \caption{English--Hungarian pairs.}\label{tab:glossary-hu}\\
  \toprule \textbf{English} & \textbf{Hungarian} \\ \midrule \endfirsthead
  \toprule \textbf{English} & \textbf{Hungarian} \\ \midrule \endhead
  \bottomrule \endlastfoot
  Students' Scientific Conference & Tudományos Diákköri Konferencia (TDK) \\
  advisor & konzulens \\
  review sheet & bírálói lap \\
  leading figure & vezérábra \\
  abstract & kivonat, összefoglaló, rezümé \\
  appendix & függelék \\
  verifiable credential (VC) & ellenőrizhető tanúsítvány \\
  verifiable data registry (VDR) & ellenőrizhető adatregiszter (also: ellenőrizhető adatnyilvántartó) \\
  self-sovereign identity (SSI) & önrendelkező identitás (also: önrendelkezésű identitás) \\
  European Digital Identity Wallet (EUDIW) & Európai Digitális Identitástárca \\
  EU Business Wallet & EU üzleti tárca \\
  permissionless & engedélyhez nem kötött \\
  self-orchestrated & önszervező \\
  self-evaluated & önkiértékelésű \\
  commitment & elköteleződés \\
  threshold-signed & küszöbaláírással ellátott \\
  race condition & versenyhelyzet \\
  check-then-act & ellenőrzést követő cselekvés \\
  liveness & elevenség \\
  adversary-driven & támadó-központú \\
  prover & bizonyító \\
  folding & hajtogatás \\
  timer event & időzítő esemény \\
  contrary-to-duty obligation & jóvátételi kötelezettség \\
  coordinated vulnerability disclosure (CVD) & koordinált sérülékenységközzététel \\
  zero-knowledge proof (ZKP) & nem-felfedő bizonyítás \\
\end{xltabular}
}
The preferred term appears first where variants are listed.
